%_____________________________________________________________________________
%=============================================================================
% data.tex v11 (24-07-2020) dibuat oleh Lionov - Informatika FTIS UNPAR
%
% Perubahan pada versi 11 (24-07-2020)
%	- Penambahan enumitem dan nosep untuk semua list, untuk menghemat kertas
%   - Bagian V: penambahan opsi Daftar Kode Program dan Daftar Notasi
%   - Bagian XIV: menjadi Bagian XVI
%   - Bagian XV: menjadi Bagian XVII
%   - Bagian XIV yang baru: untuk pilihan jenis tanda tangan mahasiswa
%   - Bagian XV yang baru: untuk pilihan munculnya tanda tangan digital untuk
%     dosen/pejabat
%_____________________________________________________________________________
%=============================================================================

%=============================================================================
% 								PETUNJUK
%=============================================================================
% Ini adalah file data (data.tex)
% Masukkan ke dalam file ini, data-data yang diperlukan oleh template ini
% Cara memasukkan data dijelaskan di setiap bagian
% Data yang WAJIB dan HARUS diisi dengan baik dan benar adalah SELURUHNYA !!
% Hilangkan tanda << dan >> jika anda menemukannya
%=============================================================================

%_____________________________________________________________________________
%=============================================================================
% 								BAGIAN 0
%=============================================================================
% Entri untuk memperbaiki posisi "DAFTAR ISI" jika tidak berada di bagian 
% tengah halaman. Sayangnya setiap sistem menghasilkan posisi yang berbeda.
% Isilah dengan 0 atau 1 (e.g. \daftarIsiError{1}). 
% Pemilihan 0 atau 1 silahkan disesuaikan dengan hasil PDF yang dihasilkan.
%=============================================================================
\daftarIsiError{0}   
%\daftarIsiError{1}   
%=============================================================================

%_____________________________________________________________________________
%=============================================================================
% 								BAGIAN I
%=============================================================================
% Tambahkan package2 lain yang anda butuhkan di sini
%=============================================================================
\usepackage{booktabs} 
\usepackage{longtable}
\usepackage{amssymb}
\usepackage{todo}
\usepackage{verbatim} 		%multiline comment
\usepackage{pgfplots}
\usepackage{enumitem}
\usepackage{amsmath}
\usepackage{subcaption}
\usepackage{algorithm}
\usepackage{algpseudocode}
\usepackage{multirow}
\usepackage{hyperref}
%\overfullrule=3mm % memperlihatkan overfull 
%=============================================================================

%_____________________________________________________________________________
%=============================================================================
% 								BAGIAN II
%=============================================================================
% Mode dokumen: menentukan halaman depan dari dokumen, apakah harus mengandung 
% prakata/pernyataan/abstrak dll (termasuk daftar gambar/tabel/isi) ?
% - final 		: hanya untuk buku skripsi, dicetak lengkap: judul ina/eng, 
%   			  pengesahan, pernyataan, abstrak ina/eng, untuk, kata 
%				  pengantar, daftar isi (daftar tabel dan gambar tetap 
%				  opsional dan dapat diatur), seluruh bab dan lampiran.
%				  Otomatis tidak ada nomor baris dan singlespacing
% - sidangakhir	: buku sidang akhir = buku final - (pengesahan + pernyataan +
%   			  untuk + kata pengantar)
%				  Otomatis ada nomor baris dan onehalfspacing 
% - sidang 		: untuk sidang 1, buku sidang = buku sidang akhir - (judul 
%				  eng + abstrak ina/eng)
%				  Otomatis ada nomor baris dan onehalfspacing
% - bimbingan	: untuk keperluan bimbingan, hanya terdapat bab dan lampiran
%   			  saja, bab dan lampiran yang hendak dicetak dapat ditentukan 
%				  sendiri (nomor baris dan spacing dapat diatur sendiri)
% Mode default adalah 'template' yang menghasilkan isian berwarna merah, 
% aktifkan salah satu mode di bawah ini :
%=============================================================================
%\mode{bimbingan} 		% untuk keperluan bimbingan
\mode{sidang} 			% untuk sidang 1
%\mode{sidangakhir} 	% untuk sidang 2 / sidang pada Skripsi 2(IF)
%\mode{final} 			% untuk mencetak buku skripsi 
%=============================================================================
%\mode{sidangakhir}

%_____________________________________________________________________________
%=============================================================================
% 								BAGIAN III
%=============================================================================
% Line numbering: penomoran setiap baris, nomor baris otomatis di-reset ke 1
% setiap berganti halaman.
% Sudah dikonfigurasi otomatis untuk mode final (tidak ada), mode sidang (ada)
% dan mode sidangakhir (ada).
% Untuk mode bimbingan, defaultnya ada (\linenumber{yes}), jika ingin 
% dihilangkan, isi dengan "no" (i.e.: \linenumber{no})
% Catatan:
% - jika nomor baris tidak kembali ke 1 di halaman berikutnya, compile kembali
%   dokumen latex anda
% - bagian ini hanya bisa diatur di mode bimbingan
%=============================================================================
%\linenumber{no} 
\linenumber{yes}
%=============================================================================

%_____________________________________________________________________________
%=============================================================================
% 								BAGIAN IV
%=============================================================================
% Linespacing: jarak antara baris 
% - single	: otomatis jika ingin mencetak buku skripsi, opsi yang 
%			     disediakan untuk bimbingan, jika pembimbing tidak keberatan 
%			     (untuk menghemat kertas)
% - onehalf	: otomatis jika ingin mencetak dokumen untuk sidang
% - double 	: jarak yang lebih lebar lagi, jika pembimbing berniat memberi 
%             catatan yg banyak di antara baris (dianjurkan untuk bimbingan)
% Catatan: bagian ini hanya bisa diatur di mode bimbingan
%=============================================================================
\linespacing{single}
%\linespacing{onehalf}
%\linespacing{double}
%=============================================================================

%_____________________________________________________________________________
%=============================================================================
% 								BAGIAN V
%=============================================================================
% Tidak semua skripsi memuat gambar, tabel, kode program, dan/atau notasi. 
% Untuk skripsi yang tidak memuat hal-hal tersebut, maka tidak diperlukan 
% Daftar Gambar, Daftar Tabel, Daftar Kode Program, dan/atau Daftar Notasi. 
% Sayangnya hal tsb sulit dilakukan secara manual karena membutuhkan 
% kedisiplinan pengguna template.  
% Jika tidak ingin menampilkan satu/lebih daftar-daftar tersebut (misalnya 
% untuk bimbingan), isi dengan "no" (e.g. \gambar{no})
%=============================================================================
\gambar{yes}
%\gambar{no}
%\tabel{yes}
\tabel{no} 
%\kode{yes}
\kode{no} 
%\notasi{yes}
\notasi{no}
%=============================================================================

%_____________________________________________________________________________
%=============================================================================
% 								BAGIAN VI
%=============================================================================
% Pada mode final, sidang da sidangkahir, seluruh bab yang ada di folder "Bab"
% dengan nama file bab1.tex, bab2.tex s.d. bab9.tex akan dicetak terurut, 
% apapun isi dari perintah \bab.
% Pada mode bimbingan, jika ingin:
% - mencetak seluruh bab, isi dengan 'all' (i.e. \bab{all})
% - mencetak beberapa bab saja, isi dengan angka, pisahkan dengan ',' 
%   dan bab akan dicetak terurut sesuai urutan penulisan (e.g. \bab{1,3,2}). 
% Catatan: Jika ingin menambahkan bab ke-3 s.d. ke-9, tambahkan file 
% bab3.tex, bab4.tex, dst di folder "Bab". Untuk bab ke-10 dan 
% seterusnya, harus dilakukan secara manual dengan mengubah file skripsi.tex
% Catatan: bagian ini hanya bisa diatur di mode bimbingan
%=============================================================================
\bab{all}
%=============================================================================

%_____________________________________________________________________________
%=============================================================================
% 								BAGIAN VII
%=============================================================================
% Pada mode final, sidang dan sidangkhir, seluruh lampiran yang ada di folder 
% "Lampiran" dengan nama file lampA.tex, lampB.tex s.d. lampJ.tex akan dicetak 
% terurut, apapun isi dari perintah \lampiran.
% Pada mode bimbingan, jika ingin:
% - mencetak seluruh lampiran, isi dengan 'all' (i.e. \lampiran{all})
% - mencetak beberapa lampiran saja, isi dengan huruf, pisahkan dengan ',' 
%   dan lampiran akan dicetak terurut sesuai urutan (e.g. \lampiran{A,E,C}). 
% - tidak mencetak lampiran apapun, isi dengan "none" (i.e. \lampiran{none})
% Catatan: Jika ingin menambahkan lampiran ke-C s.d. ke-I, tambahkan file 
% lampC.tex, lampD.tex, dst di folder Lampiran. Untuk lampiran ke-J dan 
% seterusnya, harus dilakukan secara manual dengan mengubah file skripsi.tex
% Catatan: bagian ini hanya bisa diatur di mode bimbingan
%=============================================================================
\lampiran{all}
%=============================================================================

%_____________________________________________________________________________
%=============================================================================
% 								BAGIAN VIII
%=============================================================================
% Data diri dan skripsi/tugas akhir
% - namanpm		: Nama dan NPM anda, penggunaan huruf besar untuk nama harus 
%				  benar dan gunakan 10 digit npm UNPAR, PASTIKAN BAHWA 
%				  BENAR !!! (e.g. \namanpm{Jane Doe}{1992710001}
% - judul 		: Dalam bahasa Indonesia, perhatikan penggunaan huruf besar, 
%				  judul tidak menggunakan huruf besar seluruhnya !!! 
% - tanggal 	: isi dengan {tangga}{bulan}{tahun} dalam angka numerik, 
%				  jangan menuliskan kata (e.g. AGUSTUS) dalam isian bulan.
%			  	  Tanggal ini adalah tanggal dimana anda akan melaksanakan 
%				  sidang ujian akhir skripsi/tugas akhir
% - pembimbing	: pembimbing anda, lihat daftar dosen di file dosen.tex
%				  jika pembimbing hanya 1, kosongkan parameter kedua 
%				  (e.g. \pembimbing{\JND}{} ), \JND adalah kode dosen
% - penguji 	: para penguji anda, lihat daftar dosen di file dosen.tex
%				  (e.g. \penguji{\JHD}{\JCD} )
% !!Lihat singkatan pembimbing dan penguji anda di file dosen.tex!!
% Petunjuk: hilangkan tanda << & >>, dan isi sesuai dengan data anda
%=============================================================================
\namanpm{Arlo Dante Hananvyasa}{6182201010}
\tanggal{10}{5}{2026}         %isi bulan dengan angka
\pembimbing{\HUH}{}   
\penguji{\GDK}{\RDL}
%=============================================================================

%_____________________________________________________________________________
%=============================================================================
% 								BAGIAN IX
%=============================================================================
% Judul dan title : judul bhs indonesia dan inggris
% - judulINA: judul dalam bahasa indonesia
% - judulENG: title in english
% Petunjuk: 
% - hilangkan tanda << & >>, dan isi sesuai dengan data anda
% - langsung mulai setelah '{' awal, jangan mulai menulis di baris bawahnya
% - gunakan \texorpdfstring{\\}{} untuk pindah ke baris baru
% - judul TIDAK ditulis dengan menggunakan huruf besar seluruhnya !!
%=============================================================================
\judulINA{Pembangunan Perangkat Lunak Permainan \textit{Colored Queens} dan Perbandingan Solver \textit{Backtracking} dengan \textit{Particle Swarm Optimization}}
\judulENG{Development of a Colored Queens Game Application and Comparative Analysis of Backtracking and Particle Swarm Optimization Solvers}
%_____________________________________________________________________________
%=============================================================================
% 								BAGIAN X
%=============================================================================
% Abstrak dan abstract : abstrak bhs indonesia dan inggris
% - abstrakINA: abstrak bahasa indonesia
% - abstrakENG: abstract in english 
% Petunjuk: 
% - hilangkan tanda << & >>, dan isi sesuai dengan data anda
% - langsung mulai setelah '{' awal, jangan mulai menulis di baris bawahnya
%=============================================================================
\abstrakINA{\textit{Colored Queens} adalah \textit{puzzle} logika berbasis papan yang dipopulerkan oleh platform LinkedIn, di mana pemain harus menempatkan tepat satu menteri pada setiap baris, kolom, dan sektor warna pada papan berukuran $n \times n$, tanpa ada dua menteri yang saling bertetangga dalam delapan arah. Permasalahan ini merupakan variasi dari N-Queens klasik dengan kendala tambahan berupa sektor warna ireguler dan aturan \textit{adjacency} lokal, sehingga tidak dapat diselesaikan dengan teknik konstruktif standar dan secara alami dapat dimodelkan sebagai \textit{Constraint Satisfaction Problem} (CSP).

Tugas akhir ini mengimplementasikan dan membandingkan dua pendekatan algoritmik untuk menyelesaikan \textit{Colored Queens}: \textit{backtracking} yang diperkuat dengan \textit{forward checking} dan \textit{Arc Consistency 3} (AC-3) sebagai pendekatan deterministik dan lengkap, serta \textit{Particle Swarm Optimization} (PSO) diskrit sebagai pendekatan \textit{metaheuristic}. Selain kedua solver tersebut, dikembangkan pula sebuah aplikasi permainan berbasis \textit{desktop} menggunakan Java Swing yang mengintegrasikan solver \textit{backtracking} sebagai mesin untuk fitur \textit{hint} dan \textit{auto-solve}. Pengujian dilakukan terhadap 1.502 puzzle yang mencakup ukuran papan $7 \times 7$ hingga $30 \times 30$, dengan dataset yang diperoleh dari 
\textit{playqueensgame.com} melalui \textit{web scraping}.

Solver \textit{backtracking} dengan FC+AC-3 berhasil menyelesaikan seluruh 1.500 puzzle ukuran standar ($7 \times 7$ hingga $12 \times 12$) dengan tingkat keberhasilan 100\% dan waktu eksekusi rata-rata kurang dari 2 ms. Pada papan berukuran besar ($20 \times 20$ dan $30 \times 30$), algoritma ini sebagai algoritma yang bersifat lengkap pada prinsipnya tetap akan menemukan solusi, namun waktu komputasi yang diperlukan tidak praktis, mengindikasikan bahwa \textit{overhead} propagasi AC-3 pada setiap simpul pencarian mengkompensasi manfaat pemangkasan domain pada skala tersebut. Solver PSO diskrit menunjukkan penurunan tingkat keberhasilan yang tajam seiring bertambahnya ukuran papan, dari 99,20\% pada papan $7 \times 7$ hingga hanya 5,20\% pada papan $12 \times 12$ dengan konfigurasi parameter terbaik, dan gagal sepenuhnya pada ukuran besar. Analisis sensitivitas parameter menunjukkan bahwa jumlah partikel merupakan faktor yang paling berpengaruh terhadap kualitas solusi, meskipun peningkatannya mengalami \textit{diminishing returns} pada ukuran papan yang lebih besar. Secara keseluruhan, hasil ini mengonfirmasi bahwa untuk CSP berkendala ketat seperti \textit{Colored Queens}, pendekatan berbasis propagasi kendala terbukti jauh lebih efektif dibandingkan pendekatan \textit{metaheuristic}.}
\abstrakENG{\textit{Colored Queens} is a logic-based board puzzle popularized by the LinkedIn platform, in which the player must place exactly one queen in each row, column, and colored region of an $n \times n$ board, with no two queens adjacent in any of eight directions. The problem is a variant of the classical N-Queens problem featuring irregular colored regions and a local adjacency constraint, making it unsolvable by standard constructive techniques and naturally suited for modeling as a Constraint Satisfaction Problem (CSP).

This thesis implements and compares two algorithmic approaches to solving \textit{Colored Queens}: backtracking augmented with forward checking and Arc Consistency 3 (AC-3) as a deterministic and complete approach, and discrete Particle Swarm Optimization (PSO) as a metaheuristic approach. In addition to the two solvers, a desktop application was developed using Java Swing, integrating the backtracking solver as the engine for hint and auto-solve features. Testing was conducted on 1,502 puzzles spanning board sizes from $7 \times 7$ to $30 \times 30$, with the dataset obtained from \textit{playqueensgame.com} via web scraping.

The backtracking solver with FC+AC-3 successfully solved all 1,500 standard-size puzzles ($7 \times 7$ to $12 \times 12$) with a 100\% success rate and an average execution time under 2 ms. On large boards ($20 \times 20$ and $30 \times 30$), the algorithm is theoretically guaranteed to find a solution by virtue of being a complete algorithm, but the required computation time is impractical, indicating that the AC-3 propagation overhead at each search node outweighs the benefit of domain pruning at that scale. The discrete PSO solver exhibited a sharp decline in success rate as board size increased, from 99.20\% on $7 \times 7$ boards to only 5.20\% on $12 \times 12$ under the best parameter configuration, and failed entirely on larger boards. Parameter sensitivity analysis revealed that swarm size was the most influential factor on solution quality, though its effect exhibited diminishing returns on larger board sizes. Overall, these results confirm that for tightly-constrained CSPs such as \textit{Colored Queens}, constraint propagation-based approaches are far more effective than metaheuristic ones.} 
%=============================================================================

%_____________________________________________________________________________
%=============================================================================
% 								BAGIAN XI
%=============================================================================
% Kata-kata kunci dan keywords : diletakkan di bawah abstrak (ina dan eng)
% - kunciINA: kata-kata kunci dalam bahasa indonesia
% - kunciENG: keywords in english
% Petunjuk: hilangkan tanda << & >>, dan isi sesuai dengan data anda.
%=============================================================================
\kunciINA{\textit{Colored Queens}, \textit{constraint satisfaction problem}, \textit{backtracking}, \textit{forward checking}, AC-3, \textit{particle swarm optimization}, \textit{puzzle} logika.}
\kunciENG{Colored Queens, constraint satisfaction problem, backtracking, forward checking, AC-3, particle swarm optimization, logic puzzle.}
%=============================================================================

%_____________________________________________________________________________
%=============================================================================
% 								BAGIAN XII
%=============================================================================
% Persembahan : kepada siapa anda mempersembahkan skripsi ini ...
% Petunjuk: hilangkan tanda << & >>, dan isi sesuai dengan data anda.
%=============================================================================
\untuk{Dipersembahkan untuk Tuhan, keluarga, dan teman-teman}
%=============================================================================

%_____________________________________________________________________________
%=============================================================================
% 								BAGIAN XIII
%=============================================================================
% Kata Pengantar: tempat anda menuliskan kata pengantar dan ucapan terima 
% kasih kepada yang telah membantu anda bla bla bla ....  
% Petunjuk: hilangkan tanda << & >>, dan isi sesuai dengan data anda.
%=============================================================================
\prakata{Puji syukur penulis panjatkan kepada Tuhan Yang Maha Esa, karena atas rahmat dan karunia-Nya penulis dapat menyelesaikan penyusunan tugas akhir yang berjudul ``Pembangunan Perangkat Lunak Permainan \textit{Colored Queens} dan Perbandingan Solver \textit{Backtracking} dengan \textit{Particle Swarm Optimization}``. Tugas akhir ini disusun sebagai salah satu syarat untuk memperoleh gelar sarjana di Program Studi Informatika Universitas Katolik Parahyangan.

Penyusunan tugas akhir ini tidak akan dapat diselesaikan tanpa bantuan dan dukungan dari berbagai pihak. Oleh karena itu, penulis ingin mengucapkan terima kasih yang sebesar-besarnya kepada:
	
\begin{enumerate}
	\item Kedua orang tua penulis, yang selalu memberikan dukungan, doa, dan motivasi 
	tanpa henti selama proses pengerjaan tugas akhir ini dan selama proses pembelajaran di Universitas Katolik Parahyangan.
		
	\item Bapak \HUH selaku dosen pembimbing yang telah meluangkan waktu, memberikan bimbingan, arahan, dan semangat sehingga penulis dapat menyelesaikan tugas akhir ini.
		
	\item Bapak \GDK dan $\ldots$ selaku dosen penguji yang telah memberikan kritik dan saran yang membangun demi perbaikan tugas akhir ini.
	
	\item Ibu \MTA yang telah memberikan kesempatan dan kepercayaan kepada penulis untuk menjalankan tugas sebagai pelatih ekstrakurikuler di SMA Bina Bakti.
	
	\item Bapak \RCP dan Bapak \HUH yang telah memberikan kesempatan dan kepercayaan kepada penulis untuk mengabdi sebgai asisten laboratorium.
	
	\item Ibu \VSM dan Ibu \VAN selaku dosen mata kuliah Proyek Sistem Informasi yang telah memberikan kesempatan dan kepercayaan bagi penulis untuk mengerjakan proyek SI CHIPS.
		
	\item Angelique Gabriella Halim, atas segala dukungan, kesabaran, dan semangat yang diberikan kepada penulis selama pengerjaan tugas akhir ini.
		
	\item Teman-teman penulis, terutama Renggana Santika Kurniawan dan Junita Hariyati atas kebersamaan, dukungan, dan bantuan selama persiapan sidang.
		
	\item Seluruh pihak yang telah memberikan dukungan kepada penulis, baik secara langsung maupun tidak langsung, yang tidak dapat disebutkan satu per satu.
\end{enumerate}
	
Penulis menyadari bahwa tugas akhir ini masih jauh dari sempurna. Oleh karena itu, penulis terbuka terhadap kritik dan saran yang membangun sebagai bahan perbaikan di masa mendatang. Semoga tugas akhir ini dapat bermanfaat bagi pembaca dan menjadi kontribusi yang berguna bagi penelitian selanjutnya.} 
%=============================================================================

%_____________________________________________________________________________
%=============================================================================
% 								BAGIAN XIV
%=============================================================================
% Jenis tandatangan di lembar pernyataan mahasiswa tentang plagiarisme.
% Ada 4 pilihan:
%   - digital   : diisi menggunakan digital signature (menggunakan pengolah
%                 pdf seperti Adobe Acrobat Reader DC).
%   - gambar    : diisi dengan gambar tandatangan mahasiswa (file tandatangan
%                 bertipe pdf/png/jpg). Dianjukan menggunakan warna biru.
%                 Letakkan gambar di folder "Gambar" dengan nama ttd.jpg/
%                 ttd.png/ttd.pdf (tergantung jenis file. Hapus file ttd.jpg
%                 yang digunakan sebagai contoh
%   - materai   : khusus bagi yang ingin mencetak buku dan menandatangani di 
%                 atas materai. Sama dengan pilihan ``digital'' dan dicetak.
%   - kosong    : tempat kosong ini bisa diisi dengan tanda tangan yang
%                 digambar langsung di atas pdf (fill&sign via acrobat, tanda
%                 tangan dapat dibuat dengan mouse atau stylus)
%=============================================================================
\ttd{digital}
%\ttd{gambar}
%\ttd{meterai}
%\ttd{kosong}
%=============================================================================
%\ttd{meterai}

%_____________________________________________________________________________
%=============================================================================
% 								BAGIAN XV
%=============================================================================
% Pilihan tanda tangan digital untuk dosen/pejabat:
%   - no    : pdf TIDAK dapat ditandatangani secara digital, mengakomodasi 
%             yang akan menandatangani via ``menulis'' di file pdf
%   - yes   : pdf dapat ditandatangani secara digital
% 
% PERHATIAN: perubahan ini harus ditanyakan ke kaprodi/dosen koordinator, 
% apakah harus mengisi ``no" atau ``yes". Default = no 
% Untuk mahasiswa Informatika = yes
%=============================================================================
\ttddosen{yes}
%=============================================================================
%\ttddosen{no}

%_____________________________________________________________________________
%=============================================================================
% 								BAGIAN XVI
%=============================================================================
% Tambahkan hyphen (pemenggalan kata) yang anda butuhkan di sini 
%=============================================================================
\hyphenation{ma-te-ma-ti-ka}
\hyphenation{fi-si-ka}
\hyphenation{tek-nik}
\hyphenation{in-for-ma-ti-ka}
%=============================================================================

%_____________________________________________________________________________
%=============================================================================
% 								BAGIAN XVII
%=============================================================================
% Tambahkan perintah yang anda buat sendiri di sini 
%=============================================================================
\renewcommand{\vtemplateauthor}{lionov}
\pgfplotsset{compat=newest}
\setlist{nosep}
%=============================================================================
